\chapter*{Abstract}
\addcontentsline{toc}{chapter}{Abstract}
\markboth{Abstract}{Abstract}
This report describes the design, implementation, verification and evaluation of \textbf{NOVA}, an interactive compiler-engineering platform built for the Compiler Design course. The platform contains \emph{two complete compilers} and the tooling around them.

\textbf{First}, \nova{} is a small statically typed programming language with integers, floats, booleans, strings, heap arrays, functions and recursion. Its compiler is written in C++17 and contains a hand-written lexer; a \emph{canonical LR(1)} parser generator (\LRStates{} states, \LRConflicts{} conflicts) with panic-mode error recovery; a typed AST; scoped semantic analysis with spelling suggestions; a three-address intermediate representation; control-flow, dominator, liveness, reaching-definition and pruned-SSA analyses; and an optimizer with seven passes including interprocedural inlining. Programs run on a stack-based bytecode virtual machine that supports time-travel debugging, or are compiled to \textbf{native x86-64 machine code} with a linear-scan register allocator.

\textbf{Second}, \novasql{} is a SQL compiler and in-memory engine: lexer, recursive-descent parser, binder with three-valued NULL logic, logical plan, rule-based optimizer (constant folding, predicate pushdown through joins, hash-join selection, top-N), expression bytecode and an executor.

\textbf{Third}, an \emph{assistant} turns compiler diagnostics into compiler-verified quick fixes, an iterative auto-fixer and rule-based runtime root-cause analysis (an optional LLM layer is re-checked by the compiler), a \emph{language server} exposes the compiler to any editor, and a React/TypeScript \emph{studio} with 19 pages visualises every phase from real compiler output.

\textbf{Correctness is the central theme.} Rather than trusting any single implementation, the project cross-checks independent ones: a reference tree-walking interpreter, the VM, the optimized VM and native code must agree on thousands of random programs, and \novasql{} is compared with SQLite on thousands of random queries. These differential tests found real bugs during development, which are documented in the report. In the final measurements, \FuzzPassed{} of \FuzzN{} random NOVA programs agree across three engines, \NativeDiffPass{} programs agree between native code and the VM, and \SQLFuzzPass{} of \SQLFuzzTotal{} random SQL queries agree with SQLite (the remaining \SQLFuzzSkipped{} were statements SQLite itself rejected). Native code runs about \BenchFibVspy$\times$ faster than CPython on a recursive benchmark, inlining makes a call-heavy loop \BenchInlineGain$\times$ faster, and the register allocator is measured honestly: it helps call-heavy code and does nothing for a divide-bound loop.

The report explains every component in plain language first and technical detail second, includes the real compiler output behind every example, lists the limitations explicitly, and ends with a user manual, a diagnostic catalogue, the full grammar and a glossary.

\vspace{1em}
\begin{center}\sffamily\small\color{grey}
\textbf{Keywords:} compiler design, LR(1) parsing, SSA, data-flow analysis, optimization, inlining, register allocation, x86-64, virtual machine, SQL query optimization, differential testing, fuzzing, language server
\end{center}
